%\needspace{4\baselineskip}
\color{uniblue}{\subsection{Токовая защита нулевой последовательности}}
\color{black}

\begin{enumerate}[label=\arabic{section}.\arabic{subsection}.\arabic{enumi}, labelsep=4pt, leftmargin=0pt, itemindent=57pt, itemsep=0pt, parsep=0pt] % если что можно крутить itemsep


\item
Для батареи статических конденсаторов предусматривается 2 комплекта токовой защиты нулевой 
последовательности, предназначенных для резервирования действия основных быстродействующих 
защит при внутренних земляных видах коротких замыканий. Вывод защиты производится накладкой 
\textbf{«ТЗНП-1 ВВ»}.

\item
Измерительные цепи первого комплекта ТЗНП включены на токи фаз А, В, С со стороны 
высоковольтных вводов БСК. Защита реагирует расчетное значение утроенного тока нулевой 
последовательности, определяемого по формуле:

\begin{equation}
\label{eq:3I0}
3\dot{I}_{0} = \dot{I}_{\textsl{A}} + a^2 \cdot \dot{I}_{\textsl{B}} + a \cdot \dot{I}_{\textsl{C}}
\end{equation}

\item
Второй комплект ТЗНП может быть включен на токи фаз А, В, С со стороны нейтральных выводов 
БСК или к измерительному трансформатору току, установленному в цепи заземления нейтрали БСК. 
Данный комплект обладает высокой чувствительностью к однофазным коротким замыканиям вблизи 
нейтрали. 

\item
Оба комплекта выполнены идентично и содержат по 2 ступени, выполненных на базе модулей 
\textbf{«ТЗНП»}.  

\item
Каждая из ступеней защиты может выполнена с блокировкой по второй гармонике (отношению 
уровня второй гармоники к первой). Данная блокировка реализована на блоке  «ОВ БНТ ВВ». Выводится 
данная блокировка накладкой \textbf{«Режим БНТ ТЗНП-1 ВВ»}. 

\item
Защита имеет регулируемые значения по току и времени срабатывания, задаваемые 
параметрами \textbf{«3I0ср ТЗНП-1 ВВ»} и \textbf{«Тср ТЗНП-1ВВ»} соответственно. 

\item
Функциональный блок «ТЗНП НВ» аналогичен блоку «ТЗНП ВВ» 


\begin{figure}[H]
\centering
\includegraphics[width=1\textwidth,height=1\textheight,keepaspectratio]{\fbpath/271.0241 ТЗНП/img/TZNP_VV.pdf}
\captionof{figure}{Функциональный блок ТЗНП ВВ}\label{pict:TZNP_VV}
\end{figure}


\small
\setlength{\extrarowheight}{0.05cm} % высота ячеек в таблице
\begin{longtable}
{|>{\centering\arraybackslash}m{4.5cm}|>{\centering\arraybackslash}m{4.5cm}|>{\centering\arraybackslash}m{3cm}|>{\centering\arraybackslash}m{3.5cm}|}
\caption{Перечень уставочных параметров ТЗНП\hfill\vspace{-0.5\baselineskip}}\label{ksn:tbl2} \\ 
\hhline{|-|-|-|-|}
\rowcolor{gray!30}
Наименование уставки & Диапазон & Шаг &  Значение по умолчанию \\ 
\hhline{|-|-|-|-|}
\endfirsthead

\caption*{\hspace{3pt}\emph{Продолжение таблицы \ref{ksn:tbl2}\hfill\vspace{-0.5\baselineskip}}} \\ % сделано по ГОСТ 2.105 п.6.8.7
\hhline{|-|-|-|-|}
\rowcolor{gray!30}
Наименование уставки & Диапазон & Шаг &  Значение по умолчанию \\  
\hhline{|-|-|-|-|}
\endhead
\hhline{|-|-|-|-|}

\endfoot
\endlastfoot

%%%%%%%%%%%%содержимое%%%%%%%%%%%%%%

\end{longtable}
\setlength{\extrarowheight}{0.0cm} % возврат параметра высота ячеек в таблице
\normalsize










































\end{enumerate} 